\begin{helper}[Missing witnesses for replacement-tree colors]
Fix a defective color $\delta$, put $r=D.\gammaIndex(\delta)$,
\[
 Z=V(T)\setminus\{C.\mathit{blockVertex}(r,0),
                   C.\mathit{blockVertex}(r,1)\},
 \qquad T_\delta=(R\delta).\mathit{replacementTree},
\]
and let $\varphi_\delta$ be the
\noderef{TwoColorComponentSwap} of $C.\gamma(r)$ and $\delta$ rooted at
$D.v0(\delta)$.  Every color in
\noderef{TreeUsedColors}$(\varphi_\delta,T_\delta)$ is missing under
$\varphi_\delta$ at some vertex of $Z$.
\end{helper}
\begin{proof}
Write $\gamma=C.\gamma(r)$.  Let $c$ be used on $T_\delta$, and choose a
noninitial displayed edge $h$ of $T_\delta$ colored $c$.  The exact
edge-set, retained-position, and last-edge fields of
\noderef{ScheideReplacementData} split $h$ into three possibilities: a
retained noninitial base edge of $T$, a retained block edge of $T$, or
one of the appended edges $(R\delta).f1,(R\delta).f2$.  The deleted block
starts at position $p-1\ge1$.  Consequently the original first edge is
still the first replacement edge, and injectivity of the edge sequence in
\noderef{TreeSequence} prevents it from occurring again; hence it cannot
be the witness $h$.

Suppose first that $h$ is a retained noninitial base edge.  The
\noderef{TashkinovTree} field of the original
\noderef{GlobalMaximumTriple} supplies an earlier vertex where its color
$c$ is missing.  The base-position fields of
\noderef{ScheideStructuralConfiguration} put that earlier vertex in the
base.  The deleted vertices are block vertices at positions at least
$p$, so vertex injectivity in \noderef{TreeSequence} puts the witness in
$Z$.  Moreover the base-avoidance field of
\noderef{ScheideStructuralConfiguration} gives $c\ne\gamma$.

Suppose next that $h$ is a retained block edge.  Its block-color and
base-missing-witness fields in
\noderef{ScheideStructuralConfiguration} give a base vertex $z$ at which
$c$ is missing.  Again $z\in Z$.  Since this is a retained block rather
than the deleted block, injectivity of the block-color map gives
$c\ne\gamma$.  Every displayed old tree edge has both endpoints in
$V(T)$ by \noderef{TreeSequence}, whereas $D.f0(\delta)$ has an endpoint
outside $V(T)$ by \noderef{ScheideDefectiveSetup}; hence the retained
edge is not $D.f0(\delta)$.

In both retained cases the fixed-edge clause of
\noderef{ScheideDefectiveSwitchEffect} preserves the edge color.  Its
missing-set clauses preserve the missing witness unless that witness is
$D.v0(\delta)$.  In the exceptional case, the erase-insert formula there
removes only $\gamma$, and $c\ne\gamma$, so it preserves $c$ as well.
Thus the retained-edge color $c$ is missing under $\varphi_\delta$ at the
chosen vertex of $Z$.

Finally suppose that $h$ is appended.  The endpoint and color fields of
\noderef{ScheideReplacementData} say that $h$ has old color $\delta$, has
an endpoint in $Z$, and is different from $D.f0(\delta)$.  The fixed-edge
clause of \noderef{ScheideDefectiveSwitchEffect} therefore leaves its
color equal to $\delta$.  The same theorem says that the new missing set
at $D.v0(\delta)$ is obtained by inserting $\delta$.  The exit field of
\noderef{ScheideDefectiveSetup}, followed by
\noderef{ScheideStructuralExitsInBase}, places $D.v0(\delta)$ in the
base, hence in $Z$.  Thus $D.v0(\delta)$ is the required missing witness
for the appended color.  These three exhaustive cases prove the claim.
\end{proof}
